\documentclass[10pt]{article}

% --- Page Setup ---
% \usepackage[a4paper, margin=0.5cm]{geometry}
\usepackage[a4paper, left=0.5cm, right=0.5cm, top=0.3cm, bottom=0.3cm]{geometry}
\usepackage{multicol}

% \allowdisplaybreaks

\usepackage{booktabs} % Required for professional horizontal rules (\toprule, \midrule, \bottomrule)
\usepackage{array}    % Required for better column formatting

% --- Math and Code ---
\usepackage{amsmath, amssymb}
\usepackage{listings}

% --- Spacing Adjustments ---
\usepackage{enumitem}
\usepackage{titlesec}

\usepackage{nicematrix, tikz}

\usepackage{graphicx}

\usepackage{float}

\usepackage[hidelinks]{hyperref} % hidelinks removes ugly red boxes around links

% 1. Remove space between list items
\setlist{nosep, leftmargin=*} 

% 2. Shrink space around section headers
\titlespacing*{\section}{0pt}{0pt}{0pt}
\titlespacing*{\subsection}{0pt}{1ex plus .2ex minus .2ex}{0.5ex plus .1ex}

% --- NEW: Reduce header font sizes ---
\titleformat*{\section}{\normalsize\bfseries\color{blue!75!black}}
\titleformat*{\subsection}{\small\bfseries\color{magenta}}

% 3. Add a thin line between columns for readability
\setlength{\columnseprule}{0.4pt}
\setlength{\columnsep}{1.5em}

% --- Custom Code Formatting ---
\lstset{
    basicstyle=\ttfamily\scriptsize,
    breaklines=true,
    aboveskip=2pt,
    belowskip=2pt
}

\usepackage{xcolor}
% \newcommand{\sectionbox}[1]{%
%      \colorbox{blue!75!black}{\color{white}\small\bfseries~#1~}%
%    }
%    \titleformat{\section}{}{}{0em}{\sectionbox}
%    \titleformat*{\subsection}{\normalsize\bfseries}

\raggedbottom

\setlength{\parskip}{0pt}
\setlength{\baselineskip}{0pt plus 0.2pt} 


\begin{document}

\section*{Computer Arithmetic}


Carry = 1 if '1' bit gets carried out of MSB. Indicates error in unsigned no. system.\\
V = 1 if 
\begin{itemize}
    \item \textbf{A+B}, MSB(A) = MSB(B) $\neq$ MSB(Result)
    \item \textbf{A-B}, MSB(A) $\neq$ MSB(B) and MSB(Result) $\neq$ MSB(A)
\end{itemize}
\textbf{Sign extension}: copy MSB to all the higher bits.\\
\textbf{Multi-precision arithmetic}: add the lower 32-bit words, then ADC the upper 32-bit words.\\
\textbf{Range of 2's comp} = $-2^{N-1}$ to $2^{N-1}-1$\\
\textbf{Fixed point}
\begin{itemize}
    \item Precision = value of LSB (smallest fractional value)
    \item $2^3\,2^2\,2^1\,2^0\,.\,2^{-1}\,2^{-2}\,2^{-3}\,2^{-4}$, precision is $2^{-4}$.
    \item Precision is fixed and uniform throughout the range.
    \item For a fixed number of bits, allocating more bits for the integer part (increase range) will reduce the precision.
\end{itemize}
\textbf{Floating point}
\begin{itemize}
    \item 32-bit: 1 (Sign), 8 (exponent, $E$), 23 (fraction, $F$), bias = 127
    \item 64-bit: 1 (Sign), 11 (exponent, $E$), 52 (fraction, $F$), bias = 1023
    \item Format: $(-1)^S \times (1.F)_2 \times 2^{E - \text{Bias}}$
    \item E.g. 1 00001100 010100....0 $\equiv$ $-1.3125\times 2^{-115}$
    \item Normalised mode
    \begin{itemize}
        \item Maintain a non-zero digit (1) before the radix point.
        \item Exponent range: 0000 0001 to 1111 1110
        \item Smallest magnitude: X 0000 0001 0000....0 $\equiv$ $1 \times 2^{-126}$
        \item Largest magnitude: X 1111 1110 1111 ....1 $\equiv$ $\approx 2 \times 2^{127} = 2^{128}$
        \item Underflow: numbers between $-2^{126}$ to $2^{126}$, i.e. values close to zero that cannot be represented.
    \end{itemize}
\end{itemize}
\textbf{Fixed vs floating pt (32-bit)}
\begin{itemize}
    \item Max range (fixed) = $2^{32}$ (no fractional part)
    \item Max range (floating) = $2*2^{128}$ ($-2^{128}$ to $\sim 2^{128}$)
    \item Max precision (fixed) = $2^{-32}$ (no integer part)
    \item Max precision (floating, near to 0) = less than $2^{-126}$
    \item Floating point has higher range and better precision at small numbers, but very coarse precision at the 2 ends of the range. Floating point provides uniform precision across entire range.
\end{itemize}
\textbf{More stuff}
\begin{itemize}
    \item Arithmetic between numbers with huge diff in magnitude: add/subtract numbers of similar smaller magnitude first, to allow their magnitude to be closer to the bigger magnitudes.
    \item To preserve precision AMAP, do division last, as division truncates LSB and loses bits.
\end{itemize}

\end{document}
